
In this paper, we studied feature learning in networks as a search-based optimization problem. This perspective gives us multiple advantages. It can explain classic search strategies on the basis of the exploration-exploitation trade-off. Additionally, it provides a degree of interpretability to the learned representations when applied for a prediction task. For instance, we observed that BFS can explore only limited neighborhoods. This makes BFS suitable for characterizing structural equivalences in network that rely on the immediate local structure of nodes. On the other hand, DFS can freely explore network neighborhoods which is important in discovering homophilous communities at the cost of high variance. 

Both DeepWalk and LINE can be seen as rigid search strategies over networks. DeepWalk~\cite{deepwalk} proposes search using uniform random walks. The obvious limitation with such a strategy is that it gives us no control over the explored neighborhoods. LINE~\cite{line} proposes primarily a breadth-first strategy, sampling nodes and optimizing the likelihood independently over only 1-hop and 2-hop neighbors. The effect of such an exploration is easier to characterize, but it is restrictive and provides no flexibility in exploring nodes at further depths. In contrast, the search strategy in \nodevec is both flexible and controllable exploring network neighborhoods through parameters $p$ and $q$. While these search parameters have intuitive interpretations, we obtain best results on complex networks when we can learn them directly from data. From a practical standpoint, \nodevec is scalable and robust to perturbations.

We showed how extensions of node embeddings to link prediction outperform popular heuristic scores designed specifically for this task. Our method permits additional binary operators beyond those listed in Table~\ref{tab:edgeop}. As a future work, we would like to explore the reasons behind the success of Hadamard operator over others, as well as establish interpretable equivalence notions for edges based on the search parameters. Future extensions of \nodevec could involve networks with special structure such as heterogeneous information networks, networks with explicit domain features for nodes and edges and signed-edge networks. Continuous feature representations are the backbone of many deep learning algorithms, and it would be interesting to use \nodevec representations as building blocks for end-to-end deep learning on graphs.




